\section{Outcomes, allocation estimands, and inference}
\label{sec:b-statistics}

\subsection{Outcomes}

The following variables are kept distinct: $S\in\{0,1\}$, the prespecified
binary valid success; $U\in[0,1]$, the prespecified terminal utility where
meaningful; $I\in\{0,1\}$, invalid output or scientific invalidity; the measured
resource vector $C$; and execution status with failure reason.  Utility is
defined before final evaluation, with its direction, normalization, clipping,
failure/invalid-output treatment and cross-task aggregation frozen.  Reference
anchors are not estimated from held-out controller outcomes.  Utility is not
scored only on successful or valid branches; the all-assigned-arm endpoint is
defined and validity is reported separately, and a missing score caused by
collection failure is not interchangeable with task failure.  Continuous
utility can be more informative but does not automatically reduce required
sample size; its variance, range and learning difficulty are verified in the
pilot before it is made primary.

When utility is primary, the two parallel channels are
\begin{equation}
 \begin{aligned}
 r_S(w)&=1-\E[S^0\mid W=w], & a_S(w)&=\E[S^v-S^0\mid W=w],\\
 r_U(w)&=1-\E[U^0\mid W=w], & a_U(w)&=\E[U^v-U^0\mid W=w],
 \end{aligned}
 \label{eq:b-channels}
\end{equation}
with $r_U$ a utility-deficit score and not a failure probability.  For a clean
risk-versus-effect argument the comparison is $r_S$ against $a_S$ on $S$, and
$r_U$ against $a_U$ on $U$; comparing binary failure risk with utility advantage
is useful operationally but mixes the effect distinction with a target change,
and that change is controlled and disclosed when it is made.

\subsection{Allocation estimands}

For outcome $O\in\{S,U\}$, let $d_q(w)$ indicate verification at capacity $q$
and $a_O$ the conditional advantage of verification.  The gain per eligible
checkpoint is
\begin{equation}
 G_O(d_q)=\E\big[d_q(W)\,(O^v-O^0)\mid M_v=1\big],
 \label{eq:b-gain}
\end{equation}
and the primary empirical comparison is
\begin{equation}
 D_O(q)=G_O(d_{\mathrm{adv}},q)-G_O(d_{\mathrm{risk}},q).
 \label{eq:b-primary-estimand}
\end{equation}
This differs by a factor of $q$ from a conditional selected-set mean when both
policies select exactly $q$; plots and units state which quantity is shown.
Full-population policy value with fallback is reported alongside, where
verification is unavailable, and both denominators --- common feasible support
and all eligible checkpoints --- are reported.

Capacities are $q=0.25$ primary, with $q\in\{0.05,0.10,0.50\}$ descriptive
unless the protocol explicitly allocates multiplicity.  This is a new choice,
not an inherited registration.

\paragraph{Batch allocation versus streaming deployment.}
\emph{Batch allocation} chooses the top-$q$ among the current batch using only
contexts, with an exact count budget.  \emph{Streaming deployment} uses
thresholds fixed on development or calibration data, and reports the achieved
rate and any budget overrun.  Batch evaluation does not by itself validate a
streaming rate guarantee; a hard batch limit requires its own operational
allocation rule, and the paper never labels a batch result as a streaming
guarantee.

\paragraph{Cost budgets.}
For cost budgets it is stated whether $C$ is verifier overhead or incremental
total cost.  Costs need not be positive if an intervention avoids future work;
a benefit/cost ratio requires positive costs and an appropriate single-constraint
relaxation, and otherwise an explicit constrained allocator is used.  Realized
future costs are never deployment inputs for ranking.

\subsection{Splits and inference}

The workflow is: qualification and pilot disjoint from final evaluation;
training fits representations and learners; validation selects architecture,
hyperparameters and candidate policies; calibration fixes conservative
thresholds and selection; final test evaluates frozen policies.  Splitting is at
the largest relevant dependence unit, and multiple roots, seeds, fault variants
or tasks generated from one template remain in one split.

On a final set of unseen templates, both arms' outcomes evaluate a frozen policy
using the same paired checkpoint records; independent extra selection
repetitions are not needed if test outcomes were never used for ranking or
fitting.  If a diagnostic does select by estimated test effects, it uses
separate repetitions and is labelled a diagnostic.  Inference uses paired
policy contrasts and template-aware uncertainty, and states whether it is
conditional on the trained model or includes retraining variability.  A cluster
bootstrap is asymptotic and is not a distribution-free finite-sample
certificate; when batch top-$q$ policies are evaluated, resampling respects
batches/clusters and recomputes the permitted context-only allocation.  Rare
high-value events, small $q$ and concentrated support can make effective sample
sizes much smaller than root counts; selected checkpoints and distinct templates
per budget and stratum are reported.

\subsection{Conservative control and safety language}

The planned route is the small-family procedure of \cref{sec:b-conservative}:
a correctly ranged one-sided bound with an explicit simultaneous error
allocation over the family's endpoints, evaluated on independent calibration
templates, with fallback to continue when no member qualifies.  The capacity
grid is
\begin{equation}
 q\in\{0.05,0.10,0.25,0.50\},
 \label{eq:main-capacities}
\end{equation}
with $0.25$ primary.  Reported no-injected-fault quantities are: average utility
change, invalidity, thresholded catastrophic regression if prespecified,
intervention rate, fallback rate and resource use.  State-harm fractions require
sufficient within-checkpoint information and their own identification analysis;
they are not computed from separate stochastic branch marginals as if they were
identified.
